\subsubsection{The Energy-Conjugate (Renton-Romano) Trace}

The energy-conjugate trace is defined by requiring that the beam strains 
$(\bm{\gamma}, \bm{\kappa})$ be work-conjugate to the classical resultants 
$(\bm{n}, \bm{m})$. 
That is, for any variation $\delta\hat{\bm{u}}$ within the 
Saint-Venant class:
\begin{equation}
\delta\bm{\gamma} \cdot \bm{n} + \delta\bm{\kappa} \cdot \bm{m} 
= \int_\Section \delta\bm{\varepsilon} : \bm{\sigma} \, \mathrm{d}A.
\label{eq:energy-conjugacy}
\end{equation}

We derive this trace by examining the structure of the Saint-Venant strain 
field and identifying which strain measures satisfy \eqref{eq:energy-conjugacy}.

\paragraph{Saint-Venant strain fields.}

From \eqref{eq:iesan-point-strain}, the pointwise strains in a Saint-Venant 
solution are:
\begin{equation}
    \bm{\varepsilon} = \underbrace{\begin{bmatrix}
        \FrameBasis\otimes\FrameBasis & \FrameBasis\otimes\bm{r} & \FrameBasis\times\bm{r} + \nabla\WarpTwistIesan
        & \mathbf{B}_{(b)}
    \end{bmatrix}}_{:=\mathbf{B}}
    \begin{pmatrix}
        \epsilon_0 \\ \bm{\epsilon}_1\\ \varkappa \\ \bm{\beta}
    \end{pmatrix}
\label{eq:sv-shear-strain}
\end{equation}
where the local parameters are 
$\hat{\bm{e}} = (\epsilon_0, \bm{\beta}, \hat{\varkappa}, \bm{\epsilon}_1) = \mathbf{X}^{-1}\,\bm{e}$ and 
\begin{equation*}
\mathbf{B}_{b} = (\nabla\WarpShearIesan)^T + \bm{W}_{(b)}^T
\qquad\text{with}\qquad
\bm{W}_{(b)} := -\nu\big(\mathrm{dev}(\bm{r} \otimes \bm{r}) 
- \bm{r} \otimes \CentroidLocation\big).
\end{equation*}

For a linear elastic material with Young's modulus $E$ and shear modulus $G$:
\begin{equation*}
\sigma = E\,\PointAxialStrain, \qquad \bm{\tau} = G\,\PointShearStrain.
\end{equation*}

The resultant reactions are related to the Ieşan parameters through a linear map:
\begin{equation}
\begin{pmatrix} 
N  \\ \ShearForce \\  T \\ \bm{M} 
\end{pmatrix}
= -\IesanMatrix \, \IesanVector
= -\begin{pmatrix}
EA & EA\CentroidLocation^\top & 0 & \mathbf{0}_{1\times2} \\
\mathbf{0}_{2\times1} & \mathbf{0}_{2\times2} & \mathbf{0}_{2\times1} & E\IntRGoRG \\
0 & \mathbf{0}_{1\times 2} & G\Jsv & \mathbf{0}_{1\times2} \\
EA \FrameBasis\times\CentroidLocation & \FrameBasis\times E\IntRoR & \mathbf{0}_{2\times 1} & \mathbf{0}_{2\times2} \\
\end{pmatrix}
\label{eq:resultant-map}
\end{equation}
where the \emph{resultant matrix} $\IesanMatrix \in \mathbb{R}^{6 \times 6}$ 
encodes the Saint-Venant constitutive relations  \eqref{eq:2-15} and \eqref{eq:2-18}.
$$
\IesanMatrix^{-1} =
\begin{pmatrix}
\displaystyle \frac{1}{EA} + \CentroidLocation\cdot (E\IntRGoRG)^{-1}\,\CentroidLocation
& \mathbf{0}_{1\times2}
& 0
& \CentroidLocation^\top (E\IntRGoRG)^{-1}\,\FrameBasis^\times
\\[6pt]
-(E\IntRGoRG)^{-1}\,\CentroidLocation
& \mathbf{0}_{2\times2}
& \mathbf{0}_{2\times1}
& -(E\IntRGoRG)^{-1}\,\FrameBasis^\times
\\[6pt]
0
& \mathbf{0}_{1\times2}
& \displaystyle \frac{1}{G\Jsv}
& \mathbf{0}_{1\times2}
\\[6pt]
\mathbf{0}_{2\times1}
& \bigl(E\IntRGoRG\bigr)^{-1}
& \mathbf{0}_{2\times1}
& \mathbf{0}_{2\times2}
\end{pmatrix}
$$

Thus the energetic trace is defined by the condition:
\[
\IesanMatrix \mathbf{X}^{-1} = \mathbf{X}^{-T}\underbrace{\int \mathbf{B}^T\mathbb{C}\mathbf{B}\dA}_{:=\bar{\mathbf{C}}} \mathbf{X}^{-1}
\quad\implies\quad 
\mathbf{X}^{-1} = \IesanMatrix^{-1} \bar{\mathbf{C}}^T
\]
with 
\[
\mathbb{C} := E\FrameBasis\otimes\FrameBasis + G\FrameBasis_{\alpha}\otimes\FrameBasis_{\alpha}
\]
Recall \eqref{eq:bishear-b-tensor}:
\begin{equation*}
\BishearTensorBB := \int\bm{W}_{(b)}^{\!\top}\,G\,\mathbf{B}_{(b)}\dA,
\qquad\text{and}\qquad 
\BishearTensorPB := \int\nabla\bm{\WarpShearIesan}\,G\,\mathbf{B}_{(b)}\dA
\end{equation*}

\paragraph{Work-conjugacy conditions.}

The virtual work of the section is:
\begin{equation}
\delta W = \int_\Section \delta\PointAxialStrain\,\sigma \, \mathrm{d}A 
+ \int_\Section \delta\PointShearStrain \cdot \bm{\tau} \, \mathrm{d}A.
\label{eq:virtual-work}
\end{equation}

We seek beam strains $(\epsilon, \bm{\gamma}_\perp, \varkappa, \bm{\kappa}_\perp)$ 
such that:
\begin{equation}
\delta W = \delta\epsilon\,N + \delta\bm{\gamma}_\perp \cdot \bm{F} 
+ \delta\varkappa\,T + \delta\bm{\kappa}_\perp \cdot \bm{M}.
\label{eq:conjugacy-target}
\end{equation}

We analyze each Ieşan mode separately.

\subparagraph{Extension.}

Consider $\delta a_\varepsilon \neq 0$ with all other parameters zero:
\begin{equation*}
\delta\PointAxialStrain = \delta a_\varepsilon, \qquad \delta\PointShearStrain = \mathbf{0}.
\end{equation*}
The virtual work is:
\begin{equation*}
\delta W = \delta a_\varepsilon \int_\Section \sigma \, \mathrm{d}A = \delta a_\varepsilon\,N.
\end{equation*}
Matching with \eqref{eq:conjugacy-target} requires:
\begin{equation*}
\boxed{\epsilon^{(\varepsilon)} = a_\varepsilon}
\end{equation*}

\subparagraph{Bending.}

Consider $\delta\bm{a}_\Omega \neq \mathbf{0}$ with all other parameters zero:
\begin{equation*}
\delta\PointAxialStrain = \bm{r} \cdot \delta\bm{a}_\Omega, \qquad \delta\PointShearStrain = \mathbf{0}.
\end{equation*}
The virtual work is:
\begin{equation*}
\delta W = \int_\Section (\bm{r} \cdot \delta\bm{a}_\Omega)\,\sigma \, \mathrm{d}A 
= \delta\bm{a}_\Omega \cdot \int_\Section \bm{r}\,\sigma \, \mathrm{d}A.
\end{equation*}

From the definition \eqref{eq:def-M}:
\begin{equation*}
\bm{M} = \FrameBasis \times \int_\Section \bm{r}\,\sigma \, \mathrm{d}A 
\quad\Longrightarrow\quad 
\int_\Section \bm{r}\,\sigma \, \mathrm{d}A = -\FrameBasis \times \bm{M} = \bm{M} \times \FrameBasis.
\end{equation*}
Thus:
\begin{equation*}
\delta W = \delta\bm{a}_\Omega \cdot (\bm{M} \times \FrameBasis) 
= (\delta\bm{a}_\Omega \times \bm{M}) \cdot \FrameBasis 
= -(\FrameBasis \times \delta\bm{a}_\Omega) \cdot \bm{M}.
\end{equation*}

For this to equal $\delta\bm{\kappa}_\perp \cdot \bm{M}$, we need:
\begin{equation*}
\boxed{\bm{\kappa}_\perp^{(a)} = -\FrameBasis \times \bm{a}_\Omega = -\FrameBasis^\times\,\bm{a}_\Omega}
\end{equation*}

\subparagraph{Torsion.}

Consider $\delta a_\vartheta \neq 0$ with all other parameters zero:
\begin{equation*}
\delta\PointAxialStrain = 0, \qquad 
\delta\PointShearStrain = \delta a_\vartheta\,(\FrameBasis \times \bm{r} + \nabla\WarpTwistIesan).
\end{equation*}
The virtual work is:
\begin{align*}
\delta W &= \int_\Section \delta a_\vartheta\,(\FrameBasis \times \bm{r} + \nabla\WarpTwistIesan) \cdot \bm{\tau} \, \mathrm{d}A \\
&= \delta a_\vartheta \int_\Section (\FrameBasis \times \bm{r}) \cdot \bm{\tau} \, \mathrm{d}A 
+ \delta a_\vartheta \int_\Section \nabla\WarpTwistIesan \cdot \bm{\tau} \, \mathrm{d}A.
\end{align*}

The first integral is $T$ by definition \eqref{eq:def-T}. For the second integral, 
we use the identity $\nabla \cdot (\WarpTwistIesan\,\bm{\tau}) = \nabla\WarpTwistIesan \cdot \bm{\tau} + \WarpTwistIesan\,\nabla \cdot \bm{\tau}$ 
and the divergence theorem:
\begin{equation*}
\int_\Section \nabla\WarpTwistIesan \cdot \bm{\tau} \, \mathrm{d}A 
= \oint_{\partial\Section} \WarpTwistIesan\,(\bm{\tau} \cdot \bm{\nu}) \, \mathrm{d}s 
- \int_\Section \WarpTwistIesan\,(\nabla \cdot \bm{\tau}) \, \mathrm{d}A = 0,
\end{equation*}
since $\bm{\tau} \cdot \bm{\nu} = 0$ on $\partial\Section$ (traction-free lateral surface) 
and $\nabla \cdot \bm{\tau} = 0$ (equilibrium in absence of body forces).

Thus $\delta W = \delta a_\vartheta\,T$, and matching requires:
\begin{equation}
\varkappa^{(\vartheta)} = a_\vartheta
\label{eq:conj-twist-torsion}
\end{equation}

\subparagraph{Flexure.}

This is the most involved case. Consider $\delta\bm{b}_{\Section} \neq \mathbf{0}$ with 
all other parameters zero:
\begin{align}
\delta\PointAxialStrain &= x\,(\bm{r} - \CentroidLocation) \cdot \delta\bm{b}_{\Section}, 
\label{eq:var-axial-flexure} \\
\delta\PointShearStrain &= \delta c_\vartheta\,(\FrameBasis \times \bm{r} + \nabla\WarpTwistIesan) 
+ \bm{W}_{(b)}\,\delta\bm{b}_{\Section} + (\nabla\bm{\WarpShearIesan})^\top\,\delta\bm{b}_{\Section},
\label{eq:var-shear-flexure}
\end{align}
where $\delta c_\vartheta = \CenterShearLocation \cdot \delta\bm{b}_{\Section}$.

\emph{Axial contribution.} From \eqref{eq:var-axial-flexure}:
\begin{align*}
\delta W_\sigma &= \int_\Section x\,(\bm{r} - \CentroidLocation) \cdot \delta\bm{b}_{\Section}\,\sigma \, \mathrm{d}A \\
&= x\,\delta\bm{b}_{\Section} \cdot \int_\Section (\bm{r} - \CentroidLocation)\,\sigma \, \mathrm{d}A \\
&= x\,\delta\bm{b}_{\Section} \cdot \left( \int_\Section \bm{r}\,\sigma \, \mathrm{d}A 
- \CentroidLocation \int_\Section \sigma \, \mathrm{d}A \right) \\
&= x\,\delta\bm{b}_{\Section} \cdot \left( \bm{M} \times \FrameBasis - \CentroidLocation\,N \right).
\end{align*}

Rearranging:
\begin{equation}
\delta W_\sigma = -x\,(\FrameBasis \times \delta\bm{b}_{\Section}) \cdot \bm{M} 
- x\,(\CentroidLocation \cdot \delta\bm{b}_{\Section})\,N.
\label{eq:axial-work-flexure}
\end{equation}

\emph{Shear contribution.} From \eqref{eq:var-shear-flexure}:
\begin{align}
\delta W_\tau &= \int_\Section \delta\PointShearStrain \cdot \bm{\tau} \, \mathrm{d}A \notag \\
&= \delta c_\vartheta \int_\Section (\FrameBasis \times \bm{r} + \nabla\WarpTwistIesan) \cdot \bm{\tau} \, \mathrm{d}A 
+ \delta\bm{b}_{\Section} \cdot \int_\Section \bm{W}_{(b)}^\top\,\bm{\tau} \, \mathrm{d}A \notag \\
&\quad + \delta\bm{b}_{\Section} \cdot \int_\Section \nabla\bm{\WarpShearIesan}\,\bm{\tau} \, \mathrm{d}A.
\label{eq:shear-work-flexure-1}
\end{align}

The first integral: By the same argument as for torsion (using $\bm{\tau} \cdot \bm{\nu} = 0$ 
and $\nabla \cdot \bm{\tau} = 0$), the $\nabla\WarpTwistIesan$ term vanishes:
\begin{equation*}
\int_\Section (\FrameBasis \times \bm{r} + \nabla\WarpTwistIesan) \cdot \bm{\tau} \, \mathrm{d}A = T.
\end{equation*}

The third integral: Using integration by parts on each column of $\bm{\WarpShearIesan}$:
\begin{equation*}
\int_\Section \nabla\WarpShearIesan_\alpha\,\bm{\tau} \, \mathrm{d}A 
= \oint_{\partial\Section} \WarpShearIesan_\alpha\,(\bm{\tau} \cdot \bm{\nu}) \, \mathrm{d}s 
- \int_\Section \WarpShearIesan_\alpha\,(\nabla \cdot \bm{\tau}) \, \mathrm{d}A = \mathbf{0}.
\end{equation*}
Thus $\int_\Section \nabla\bm{\WarpShearIesan}\,\bm{\tau} \, \mathrm{d}A = \mathbf{0}$.

Substituting back:
\begin{equation}
\delta W_\tau = (\CenterShearLocation \cdot \delta\bm{b}_{\Section})\,T 
+ \delta\bm{b}_{\Section} \cdot \int_\Section \bm{W}_{(b)}^\top\,\bm{\tau} \, \mathrm{d}A.
\label{eq:shear-work-flexure-2}
\end{equation}

\emph{Total virtual work from flexure.} Combining \eqref{eq:axial-work-flexure} 
and \eqref{eq:shear-work-flexure-2}:
\begin{align}
\delta W^{(b)} &= -x\,(\CentroidLocation \cdot \delta\bm{b}_{\Section})\,N 
- x\,(\FrameBasis \times \delta\bm{b}_{\Section}) \cdot \bm{M} \notag \\
&\quad + (\CenterShearLocation \cdot \delta\bm{b}_{\Section})\,T 
+ \delta\bm{b}_{\Section} \cdot \int_\Section \bm{W}_{(b)}^\top\,\bm{\tau} \, \mathrm{d}A.
\label{eq:total-work-flexure}
\end{align}

\emph{Matching with conjugacy condition.} We require:
\begin{equation}
\delta W^{(b)} = \delta\epsilon^{(b)}\,N + \delta\bm{\gamma}_\perp^{(b)} \cdot \bm{F} 
+ \delta\varkappa^{(b)}\,T + \delta\bm{\kappa}_\perp^{(b)} \cdot \bm{M}.
\label{eq:match-flexure}
\end{equation}

Comparing term by term:
\begin{itemize}
\item Coefficient of $N$: $\delta\epsilon^{(b)} = -x\,(\CentroidLocation \cdot \delta\bm{b}_{\Section})$
\item Coefficient of $\bm{M}$: $\delta\bm{\kappa}_\perp^{(b)} = -x\,(\FrameBasis \times \delta\bm{b}_{\Section})$
\item Coefficient of $T$: $\delta\varkappa^{(b)} = \CenterShearLocation \cdot \delta\bm{b}_{\Section}$
\item Coefficient of $\bm{F}$: This requires careful treatment of the remaining term.
\end{itemize}

The remaining term is $\delta\bm{b}_{\Section} \cdot \int_\Section \bm{W}_{(b)}^\top\,\bm{\tau} \, \mathrm{d}A$. 
For this to equal $\delta\bm{\gamma}_\perp^{(b)} \cdot \bm{F}$, we need to express 
it in terms of the shear force $\bm{F} = \int_\Section \bm{\tau} \, \mathrm{d}A$.

Define the \emph{Poisson-shear tensor}:
\begin{equation}
\mathbf{P} := \int_\Section \bm{W}_{(b)}^\top\,\bm{\tau} \, \mathrm{d}A 
\quad \in \mathbb{R}^2,
\label{eq:def-P}
\end{equation}
which depends on the stress state $\bm{\tau}$. Under flexure loading, 
$\bm{\tau} = G\,\PointShearStrain^{(b)}$ where:
\begin{equation*}
\PointShearStrain^{(b)} = c_\vartheta\,(\FrameBasis \times \bm{r} + \nabla\WarpTwistIesan) 
+ \bm{W}_{(b)}\,\bm{b}_{\Section} + (\nabla\bm{\WarpShearIesan})^\top\,\bm{b}_{\Section}.
\end{equation*}

The shear force under flexure is:
\begin{align}
\bm{F} &= G \int_\Section \PointShearStrain^{(b)} \, \mathrm{d}A \notag \\
&= G\,c_\vartheta \int_\Section (\FrameBasis \times \bm{r} + \nabla\WarpTwistIesan) \, \mathrm{d}A 
+ G \int_\Section \bm{W}_{(b)} \, \mathrm{d}A\,\bm{b}_{\Section} 
+ G \int_\Section (\nabla\bm{\WarpShearIesan})^\top \, \mathrm{d}A\,\bm{b}_{\Section}.
\label{eq:F-flexure}
\end{align}

Using $\int_\Section (\FrameBasis \times \bm{r}) \, \mathrm{d}A = A\,(\FrameBasis \times \CentroidLocation)$ 
and noting that $\int_\Section \nabla\WarpTwistIesan \, \mathrm{d}A = \oint_{\partial\Section} \WarpTwistIesan\,\bm{\nu} \, \mathrm{d}s$ 
(which need not vanish), define:
\begin{align}
\bm{q}_\varphi &:= \int_\Section (\FrameBasis \times \bm{r} + \nabla\WarpTwistIesan) \, \mathrm{d}A, 
\label{eq:def-q-phi} \\
\mathbf{Q}_W &:= \int_\Section \bm{W}_{(b)} \, \mathrm{d}A = A\,\SectionAverage{\bm{W}_{(b)}}, 
\label{eq:def-Q-W} \\
\mathbf{Q}_\psi &:= \int_\Section (\nabla\bm{\WarpShearIesan})^\top \, \mathrm{d}A.
\label{eq:def-Q-psi}
\end{align}

Then:
\begin{equation}
\bm{F} = G\,c_\vartheta\,\bm{q}_\varphi + G\,(\mathbf{Q}_W + \mathbf{Q}_\psi)\,\bm{b}_{\Section} 
= G\,(\CenterShearLocation \cdot \bm{b}_{\Section})\,\bm{q}_\varphi + G\,\mathbf{Q}\,\bm{b}_{\Section},
\label{eq:F-flexure-compact}
\end{equation}
where $\mathbf{Q} := \mathbf{Q}_W + \mathbf{Q}_\psi \in \mathbb{R}^{3 \times 2}$.

\emph{The shear conjugacy problem.} We need $\delta\bm{\gamma}_\perp^{(b)} \cdot \bm{F}$ 
to equal $\delta\bm{b}_{\Section} \cdot \mathbf{P}$. 
This is a compatibility condition 
between the shear strain definition and the Poisson-shear coupling.

In general, $\mathbf{P}$ and $\bm{F}$ are not simply related by a constant tensor 
acting on $\delta\bm{b}_{\Section}$. However, we can define the energy-conjugate shear 
strain implicitly.

Write:
\begin{equation}
\delta\bm{b}_{\Section} \cdot \mathbf{P} = \delta\bm{b}_{\Section} \cdot \mathbf{P}(\bm{b}_{\Section}),
\end{equation}
where $\mathbf{P}(\bm{b}_{\Section})$ is linear in $\bm{b}_{\Section}$ through the stress $\bm{\tau}$.

For the constitutive relation to be symmetric, we require:
\begin{equation}
\bm{\gamma}_\perp^{(b)} = \mathbf{K}^{-1}\,\bm{F},
\label{eq:romano-shear-def}
\end{equation}
where $\mathbf{K}$ is the shear stiffness tensor defined by $\bm{F} = \mathbf{K}\,\bm{\gamma}_\perp$.

From \eqref{eq:F-flexure-compact}, the shear force is linear in $\bm{b}_{\Section}$:
\begin{equation}
\bm{F} = \mathbf{R}\,\bm{b}_{\Section}, \qquad 
\mathbf{R} := G\,(\bm{q}_\varphi \otimes \CenterShearLocation + \mathbf{Q}).
\label{eq:F-R-relation}
\end{equation}

For energy conjugacy with symmetric stiffness, Romano defines:
\begin{equation}
\bm{\gamma}_\perp^{(b)} := \mathbf{K}^{-1}\,\mathbf{R}\,\bm{b}_{\Section},
\label{eq:ec-shear-final}
\end{equation}
where the shear stiffness $\mathbf{K}$ is determined by the requirement that 
the strain energy $\frac{1}{2}\bm{\gamma}_\perp \cdot \bm{F}$ equals the actual 
shear strain energy $\frac{1}{2}\int_\Section \PointShearStrain \cdot \bm{\tau} \, \mathrm{d}A$.

This gives:
\begin{equation}
\mathbf{K} = \mathbf{R}\,\mathbf{H}^{-1}\,\mathbf{R}^\top,
\label{eq:K-formula}
\end{equation}
where $\mathbf{H}$ is the shear energy tensor:
\begin{equation}
\mathbf{H} := G \int_\Section \PointShearStrain^{(b)} \otimes \PointShearStrain^{(b)} \, \mathrm{d}A 
\bigg|_{\bm{b}_{\Section} = \text{unit}}.
\label{eq:def-H}
\end{equation}

Substituting \eqref{eq:K-formula} into \eqref{eq:ec-shear-final}:
\begin{equation}
\boxed{\bm{\gamma}_\perp^{(b)} = \mathbf{H}^{-1}\,\mathbf{R}^\top\,(\mathbf{R}\,\mathbf{H}^{-1}\,\mathbf{R}^\top)^{-1}\,\mathbf{R}\,\bm{b}_{\Section} 
= \mathbf{H}^{-1}\,\mathbf{R}^\top\,\bm{\lambda}}
\label{eq:ec-shear-explicit}
\end{equation}
where $\bm{\lambda}$ is a Lagrange multiplier enforcing the force-strain relation.

For the simpler case where we directly define shear strain from the force:
\begin{equation}
\boxed{\bm{\gamma}_\perp^{(b)} = \mathbf{K}^{-1}\,\bm{F} = \mathbf{K}^{-1}\,\mathbf{R}\,\bm{b}_{\Section}}
\end{equation}

\paragraph{Summary of energy-conjugate strains.}

Collecting results from all modes:
\begin{align}
\epsilon &= a_\varepsilon - x\,(\CentroidLocation \cdot \bm{b}_{\Section}), 
\label{eq:ec-axial-final} \\
\bm{\gamma}_\perp &= \mathbf{K}^{-1}\,\mathbf{R}\,\bm{b}_{\Section}, 
\label{eq:ec-shear-final2} \\
\varkappa &= a_\vartheta + \CenterShearLocation \cdot \bm{b}_{\Section}, 
\label{eq:ec-twist-final} \\
\bm{\kappa}_\perp &= -\FrameBasis^\times\,\bm{a}_\Omega - x\,(\FrameBasis \times \bm{b}_{\Section}).
\label{eq:ec-curv-final}
\end{align}

\paragraph{Assembly of trace matrices.}

The beam strains $\bm{e}(x) = (\epsilon, \bm{\gamma}_\perp, \varkappa, \bm{\kappa}_\perp)^\top$ 
satisfy:
\begin{equation}
\bm{e}(x) = \TraceDeDpRoot\,\bm{p} + x\,\TraceDeDpOne\,\bm{p},
\end{equation}
with $\bm{p} = (a_\varepsilon, \bm{a}_\Omega, a_\vartheta, \bm{b}_{\Section})^\top$ and:

\begin{equation}
\TraceDeDpRoot = \begin{pmatrix}
1 & \mathbf{0}^\top & 0 & \mathbf{0}^\top \\[4pt]
\mathbf{0} & \mathbf{0} & \mathbf{0} & \mathbf{K}^{-1}\mathbf{R} \\[4pt]
0 & \mathbf{0}^\top & 1 & \CenterShearLocation^\top \\[4pt]
\mathbf{0} & -\FrameBasis^\times & \mathbf{0} & \mathbf{0}
\end{pmatrix}
\label{eq:T0-romano-final}
\end{equation}

\begin{equation}
\TraceDeDpOne = \begin{pmatrix}
0 & \mathbf{0}^\top & 0 & -\CentroidLocation^\top \\[4pt]
\mathbf{0} & \mathbf{0} & \mathbf{0} & \mathbf{0} \\[4pt]
0 & \mathbf{0}^\top & 0 & \mathbf{0}^\top \\[4pt]
\mathbf{0} & \mathbf{0} & \mathbf{0} & -\FrameBasis^\times
\end{pmatrix}
\label{eq:T1-romano-final}
\end{equation}

\paragraph{Warping index.}

The stacked matrix is:
\begin{equation}
\mathbf{T}_{\rm stack} = \begin{pmatrix} \TraceDeDpRoot \\ \TraceDeDpOne \end{pmatrix} 
\in \mathbb{R}^{12 \times 6}.
\end{equation}

Column analysis:
\begin{itemize}
\item $a_\varepsilon$: Column $(1, \mathbf{0}, 0, \mathbf{0}; 0, \mathbf{0}, 0, \mathbf{0})^\top$ — rank 1.
\item $\bm{a}_\Omega$: Columns with $-\FrameBasis^\times$ in curvature block — rank 2.
\item $a_\vartheta$: Column $(0, \mathbf{0}, 1, \mathbf{0}; 0, \mathbf{0}, 0, \mathbf{0})^\top$ — rank 1.
\item $\bm{b}_{\Section}$: In $\TraceDeDpRoot$: $\mathbf{K}^{-1}\mathbf{R}$ (shear) and $\CenterShearLocation^\top$ (twist). 
In $\TraceDeDpOne$: $-\CentroidLocation^\top$ (axial) and $-\FrameBasis^\times$ (curvature).
\end{itemize}

The $\bm{b}_{\Section}$ columns contribute rank 2 through $\TraceDeDpOne$ (via $-\FrameBasis^\times$), 
independent of the contributions from $\bm{a}_\Omega$ (which appear in different rows 
of the stacked matrix).

\begin{proposition}[Warping index of energy-conjugate trace]
\label{prop:warp-index-romano-final}
For sections with non-degenerate geometry, the energy-conjugate trace has 
warping index:
\begin{equation}
w = 6 - \mathrm{rank}(\mathbf{T}_{\rm stack}) = 0.
\end{equation}
\end{proposition}

\begin{proof}
The six columns of $\mathbf{T}_{\rm stack}$ are:
\begin{enumerate}
\item $a_\varepsilon$: Nonzero in row 1 of $\TraceDeDpRoot$.
\item $a_{\Omega,1}$: Nonzero in rows 5--6 of $\TraceDeDpRoot$ (via $-\FrameBasis^\times$).
\item $a_{\Omega,2}$: Nonzero in rows 5--6 of $\TraceDeDpRoot$ (via $-\FrameBasis^\times$).
\item $a_\vartheta$: Nonzero in row 3 of $\TraceDeDpRoot$.
\item $b_{\Omega,1}$: Nonzero in row 1 of $\TraceDeDpOne$ (via $-\CentroidLocation^\top$) 
and rows 11--12 of $\TraceDeDpOne$ (via $-\FrameBasis^\times$).
\item $b_{\Omega,2}$: Same structure as $b_{\Omega,1}$.
\end{enumerate}
These columns are linearly independent for generic $\CentroidLocation \neq \mathbf{0}$ 
(non-centroidal origin), giving $\mathrm{rank}(\mathbf{T}_{\rm stack}) = 6$.

For centroidal coordinates ($\CentroidLocation = \mathbf{0}$), the $\bm{b}_{\Section}$ columns 
still contribute through $-\FrameBasis^\times$ in $\TraceDeDpOne$ and $\mathbf{K}^{-1}\mathbf{R}$ 
in $\TraceDeDpRoot$, maintaining full rank.
\end{proof}

\begin{remark}[Constitutive symmetry]
By construction, the energy-conjugate trace yields a symmetric section 
stiffness matrix:
\begin{equation}
\bm{s} = \mathbf{C}\,\bm{e}, \qquad \mathbf{C} = \mathbf{C}^\top,
\end{equation}
where $\bm{s} = (N, \bm{F}, T, \bm{M})^\top$. This follows from the variational 
derivation: the strains are defined as derivatives of the complementary 
energy with respect to the resultants.
\end{remark}

\begin{remark}[Comparison with section-average trace]
The two traces yield identical results for:
\begin{itemize}
\item Axial strain from extension: $\epsilon^{(\varepsilon)} = a_\varepsilon$
\item Curvature from bending: $\bm{\kappa}_\perp^{(a)} = -\FrameBasis^\times\,\bm{a}_\Omega$
\item Twist from torsion: $\varkappa^{(\vartheta)} = a_\vartheta$
\item Curvature from flexure: $\bm{\kappa}_\perp^{(b)} = -x\,(\FrameBasis \times \bm{b}_{\Section})$
\end{itemize}
They differ in:
\begin{itemize}
\item Shear strain from flexure: section-average gives $\bm{\Gamma}_0\,\bm{b}_{\Section}$; 
energy-conjugate gives $\mathbf{K}^{-1}\mathbf{R}\,\bm{b}_{\Section}$.
\item Twist from flexure: section-average includes additional coupling terms; 
energy-conjugate gives $\CenterShearLocation \cdot \bm{b}_{\Section}$.
\end{itemize}
The difference is a gauge transformation that preserves the warping index 
but changes the constitutive coupling structure.
\end{remark}

\begin{remark}[Plane-section property]
The energy-conjugate trace does not satisfy the plane-section property. 
For a section-rigid displacement $\hat{\bm{u}} = \bm{a} + \bm{\omega} \times \bm{r}$, 
the extracted beam strains $(\bm{\gamma}, \bm{\kappa})$ are not generally 
$(\bm{a}' + \FrameBasis \times \bm{\omega}, \bm{\omega}')$ evaluated at a consistent 
reference point. 
\end{remark}